\documentclass[a4paper,12pt]{article}

% 页面设置
\usepackage[top=2.54cm, bottom=2.54cm, left=1.91cm, right=1.91cm]{geometry}

% 中文支持
\usepackage[UTF8]{ctex}
\usepackage{fontspec}

% 数学公式支持
\usepackage{amsmath}
\usepackage{amssymb}
\usepackage{amsfonts}
\usepackage{amsthm}

\usepackage{enumitem}

% 定理环境定义
\newtheorem*{pf}{Proof}
\newtheorem*{sol}{Solution}

% 图形支持
\usepackage{graphicx}
\usepackage{tikz}

% 超链接支持
\usepackage{hyperref}
\hypersetup{
    colorlinks=true,
    linkcolor=blue,
    filecolor=magenta,
    urlcolor=cyan,
}

% 行距设置
\linespread{1.25}

% 标题信息
\title{Mathematical Analysis II\\Week 3 Homework (March 19)}
\author{袁子强\hspace{1cm} 2025533009}
% \date{}

\begin{document}

\maketitle

\subsection*{Section 9.1}

17. 研究下列函数的连续性:

(1) $f(x,y) =
    \begin{cases}
        \dfrac{xy}{x-y}, & x \neq y, \\
        0,               & x = y;
    \end{cases}$
\begin{sol}
    令 $y=x-x^2$，则 $x-y=x^2$，于是
    $$f(x,x-x^2)=\frac{x^2-x^3}{x^2}=1-x.$$
    因此
    $$\lim_{x\to 0}f(x,x-x^2)=1-0=1\neq f(0,0)=0.$$
    故 $f(x,y)$ 在 $x=y$ 上不连续，在其他位置连续（因为当 $x\neq y$ 时，$f(x,y)$ 为初等函数）。
\end{sol}

(2) $f(x,y) =
    \begin{cases}
        x \sin \dfrac{1}{y}, & y \neq 0, \\
        0,                   & y = 0;
    \end{cases}$
\begin{sol}
    当 $y\neq 0$ 时，$f(x,y)$ 为初等连续函数。

    由于 $-1\leq\sin\dfrac{1}{y}\leq 1$，因此 $|f(x,y)|\leq|x|$。从而
    $$\lim_{x\to 0}|f(x,y)|\leq\lim_{x\to 0}|x|=0.$$
    故 $\lim\limits_{x\to 0,y\to 0}f(x,y)=f(0,0)=0$，即 $f$ 在 $(0,0)$ 处连续。

    取 $x=x_0\neq 0$，取 $y_n=\dfrac{1}{2n\pi+\dfrac{\pi}{2}}$，则当 $n\to\infty$ 时 $y_n\to 0$。此时
    $$\sin\frac{1}{y_n}=\sin\left(2n\pi+\frac{\pi}{2}\right)=1,$$
    $$f\left(x_0,\frac{1}{2n\pi+\frac{\pi}{2}}\right)=x_0\neq 0.$$
    因此 $\lim\limits_{y\to 0}f\left(x_0,\dfrac{1}{2n\pi+\dfrac{\pi}{2}}\right)=x_0\neq 0=f(x_0,0)$。
    故当 $x\neq 0,y=0$ 时，$f$ 不连续。
\end{sol}

(3) $f(x,y) =
    \begin{cases}
        \dfrac{x^2 y}{x^2 + y^2}, & (x,y) \neq (0,0), \\
        0,                        & (x,y) = (0,0);
    \end{cases}$
\begin{sol}
    当 $(x,y)\neq(0,0)$ 时，$f$ 为连续初等函数。

    考虑极坐标变换 $x=r\cos\theta$, $y=r\sin\theta$，则
    $$f(x,y)=f(r\cos\theta,r\sin\theta)=\frac{r^3\cos^2\theta\sin\theta}{r^2}=r\cos^2\theta\sin\theta.$$
    由于 $|\cos\theta|,|\sin\theta|\leq 1$，因此 $|\cos^2\theta\sin\theta|\leq 1$，从而
    $$|r\cos^2\theta\sin\theta|\leq|r|.$$
    故
    $$\lim_{r\to 0}|f(r\cos\theta,r\sin\theta)|\leq\lim_{r\to 0}|r|=0=f(0,0).$$
    因此 $f$ 在 $(0,0)$ 处连续，从而 $f$ 在 $\mathbb{R}^2$ 上处处连续。
\end{sol}

(4) $f(x,y) =
    \begin{cases}
        \dfrac{x-y}{x+y}, & x + y \neq 0, \\
        0,                & x + y = 0.
    \end{cases}$
\begin{sol}
    当 $x + y \neq 0$ 时，$f$ 为连续初等函数。

    取 $y=4x$，此时 $y+x=0$ 当且仅当 $y=x=0$。则
    $$\lim_{x\to 0}\frac{x-4x}{x+4x}=\frac{-3}{5}\neq 0=f(0,0).$$
    故 $f$ 在 $x+y=0$ 上处处不连续。
\end{sol}


18. 证明函数 $f(x,y) =
    \begin{cases}
        \dfrac{x^2 y}{x^4 + y^2}, & x^2 + y^2 > 0, \\
        0,                        & x^2 + y^2 = 0
    \end{cases}$
在点 $(0,0)$ 沿着过此点的每一射线 $x = t \cos \alpha,\ y = t \sin \alpha,\ (0 \leqslant t < +\infty)$ 连续，即 $\lim_{t \to 0} f(t \cos \alpha, t \sin \alpha) = f(0,0)$. 但此函数在点 $(0,0)$ 并不连续.
\begin{pf}
    当 $t\neq 0$ 时，
    $$f(t \cos \alpha, t \sin \alpha)=\dfrac{t^3\cos^2\alpha\sin\alpha}{t^4\cos^4\alpha + t^2\sin^2\alpha}=\frac{\cos^2\alpha\sin\alpha}{t\cos^4\alpha+\dfrac{\sin^2\alpha}{t}}.$$

    当 $\sin\alpha=0$ 时，$f(t\cos\alpha,0)=\dfrac{0}{t\cos^4\alpha}=0$，因此 $\lim\limits_{t\to 0}f(t\cos\alpha,0)=0=f(0,0)$.

    当 $\sin\alpha\neq 0$ 时，由于 $\lim\limits_{t\to 0}t\cos^4\alpha=0$，故
    $$\lim_{t\to 0}f(t \cos \alpha, t \sin \alpha)=\lim_{t\to 0}\frac{t\cos^2\alpha}{\sin\alpha}=0=f(0,0).$$

    综上所述，$\lim\limits_{t \to 0} f(t \cos \alpha, t \sin \alpha) = f(0,0)$.

    然而，令 $y=x^2$，则
    $$f(x,x^2)=\frac{x^4}{2x^4}=\frac{1}{2},$$
    从而 $\lim\limits_{x\to 0}f(x,x^2)=\dfrac{1}{2}\neq 0=f(0,0)$，故函数在点 $(0,0)$ 不连续。

    Q.E.D.
\end{pf}

19. 设 $f(x,y)$ 在 $D \subset \mathbb{R}^2$ 上分别对 $x$ 和 $y$ 连续，且关于变量 $y$ 是单调的，证明：$f(x,y)$ 在 $D$ 上连续.
\begin{pf}
    由于 $f(x,y)$ 在 $D \subset \mathbb{R}^2$ 上分别对 $x$ 和 $y$ 连续，任取 $P(x_0,y_0)\in D$，对任意 $\varepsilon>0$，存在 $\delta_2=\delta_2(\varepsilon)>0$，使得对任意 $|y-y_0|\leq\delta_2$，有
    $$|f(x_0,y)-f(x_0,y_0)|<\varepsilon.$$
    特别地，
    $$|f(x_0,y_0-\delta_2)-f(x_0,y_0)|<\varepsilon,\quad |f(x_0,y_0+\delta_2)-f(x_0,y_0)|<\varepsilon.$$

    存在 $\delta_1>0$，使得对任意 $|x-x_0|<\delta_1$，有
    $$|f(x,y_0-\delta_2)-f(x_0,y_0-\delta_2)|<\frac{\varepsilon}{2},\quad |f(x,y_0+\delta_2)-f(x_0,y_0+\delta_2)|<\frac{\varepsilon}{2}.$$

    不妨设 $f$ 关于 $y$ 单调递增，则当 $|x-x_0|<\delta_1$，$|y-y_0|<\delta_2$ 时，有
    $$f(x,y_0-\delta_2)\leq f(x,y)\leq f(x,y_0+\delta_2).$$

    因此
    $$
        \begin{aligned}
            f(x,y)-f(x_0,y_0) & \leq f(x,y_0+\delta_2)-f(x_0,y_0)                                      \\
                              & = f(x,y_0+\delta_2)-f(x_0,y_0+\delta_2)+f(x_0,y_0+\delta_2)-f(x_0,y_0) \\
                              & < \frac{\varepsilon}{2}+\frac{\varepsilon}{2}=\varepsilon,
        \end{aligned}
    $$
    $$
        \begin{aligned}
            f(x,y)-f(x_0,y_0) & \geq f(x,y_0-\delta_2)-f(x_0,y_0)                                      \\
                              & = f(x,y_0-\delta_2)-f(x_0,y_0-\delta_2)+f(x_0,y_0-\delta_2)-f(x_0,y_0) \\
                              & > -\frac{\varepsilon}{2}-\frac{\varepsilon}{2}=-\varepsilon.
        \end{aligned}
    $$

    故 $|f(x,y)-f(x_0,y_0)|<\varepsilon$，因此 $f$ 在 $P$ 处连续，从而 $f$ 在 $D$ 上连续。

    Q.E.D.
\end{pf}


20. 设 $D \subset \mathbb{R}^2$ 对任意 $(x,y) \in D$, 令 $f(x,y) = x$. 称其为 $D$ 在 $x$ 轴的投影函数. 证明投影函数是连续函数，但是它不一定将闭集映成闭集.
\begin{pf}
    任取 $P(x_0,y_0)\in D$，对任意 $\varepsilon>0$，取 $\delta=\varepsilon$，则对任意满足 $\sqrt{(x-x_0)^2+(y-y_0)^2}<\delta$ 的点 $(x,y)$，有
    $$|f(x,y)-f(x_0,y_0)|=|x-x_0|\leq\sqrt{(x-x_0)^2+(y-y_0)^2}<\delta=\varepsilon.$$
    故投影函数在 $D$ 上连续。

    考虑 $D: xy=1$，即 $xy-1=0$。令 $g(x,y)=xy-1$，由于 $g$ 在 $\mathbb{R}^2$ 上连续，且 $D=g^{-1}(\{0\})$，故 $xy=1$ 为闭集。在此情况下，$f(D)=(-\infty,0)\cup(0,+\infty)$ 为开区间，不是闭集。因此投影函数不一定将闭集映成闭集。

    Q.E.D.
\end{pf}


22. 设 $f(x,y)$ 在 $(x_0,y_0)$ 处连续，$x = x(u,v),\ y = y(u,v)$ 在 $(u_0,v_0)$ 处连续。用 $\varepsilon-\delta$ 语言证明复合函数 $f(x(u,v), y(u,v))$ 在 $(u_0,v_0)$ 处连续.
\begin{pf}
    由于 $f$ 在 $(x_0,y_0)$ 处连续，故对任意 $\varepsilon>0$，存在 $\delta_1=\delta_1(\varepsilon)>0$，使得对任意满足 $\sqrt{(x-x_0)^2+(y-y_0)^2}<\delta_1$ 的点 $(x,y)$，有
    $$|f(x,y)-f(x_0,y_0)|<\varepsilon.$$

    又由于 $x(u,v)$ 和 $y(u,v)$ 均在 $(u_0,v_0)$ 处连续，故对上述 $\delta_1>0$，存在 $\delta_2=\delta_2(\delta_1)>0$ 和 $\delta_3=\delta_3(\delta_1)>0$，使得：

    对任意满足 $\sqrt{(u-u_0)^2+(v-v_0)^2}<\delta_2$ 的点 $(u,v)$，有
    $$|x(u,v)-x(u_0,v_0)|<\frac{\delta_1}{\sqrt{2}};$$

    对任意满足 $\sqrt{(u-u_0)^2+(v-v_0)^2}<\delta_3$ 的点 $(u,v)$，有
    $$|y(u,v)-y(u_0,v_0)|<\frac{\delta_1}{\sqrt{2}}.$$

    因此，对任意 $\varepsilon>0$，取 $\delta=\min\{\delta_2,\delta_3\}$，则当 $\sqrt{(u-u_0)^2+(v-v_0)^2}<\delta$ 时，有
    $$|x(u,v)-x(u_0,v_0)|<\frac{\delta_1}{\sqrt{2}},\quad |y(u,v)-y(u_0,v_0)|<\frac{\delta_1}{\sqrt{2}},$$
    从而
    $$\sqrt{(x-x_0)^2+(y-y_0)^2}<\delta_1,$$
    故 $|f(x(u,v),y(u,v))-f(x_0,y_0)|<\varepsilon$.

    因此 $f(x(u,v), y(u,v))$ 在 $(u_0,v_0)$ 处连续。

    Q.E.D.
\end{pf}


23. 设 $f(x,y) = \dfrac{1}{1-xy},\quad (x,y) \in [0,1] \times [0,1],\ (x,y) \neq (1,1)$，证明函数连续但不一致连续.
\begin{pf}
    由于 $f$ 在定义域 $D=[0,1]\times[0,1]\setminus\{(1,1)\}$ 内有定义，且为初等函数，故 $f$ 在 $D$ 内连续。

    下证 $f$ 不一致连续。取 $\varepsilon_0=1>0$，对任意 $\delta>0$，取正整数 $n$ 使得 $0<\dfrac{1}{n}<\dfrac{\delta}{2}$ 且 $\dfrac{1}{n}<1$。

    令 $P\left(1-\dfrac{1}{n},1-\dfrac{1}{n}\right)$，$Q\left(1-\dfrac{1}{2n},1-\dfrac{1}{n}\right)$，由于 $0\leq 1-\dfrac{1}{n},1-\dfrac{1}{2n}<1$，故 $P,Q$ 均在定义域 $D$ 内。二者的距离为
    $$d(P,Q)=\sqrt{\left[\left(1-\frac{1}{2n}\right)-\left(1-\frac{1}{n}\right)\right]^2+\left[\left(1-\frac{1}{n}\right)-\left(1-\frac{1}{n}\right)\right]^2}=\frac{1}{2n}<\frac{\delta}{2}<\delta.$$

    计算得
    $$f(P)=\frac{1}{1-\left(1-\frac{1}{n}\right)^2}=\frac{1}{\frac{2}{n}-\frac{1}{n^2}}=\frac{n^2}{2n-1},$$
    $$f(Q)=\frac{1}{1-\left(1-\frac{1}{2n}\right)\left(1-\frac{1}{n}\right)}=\frac{1}{\frac{3}{2n}-\frac{1}{2n^2}}=\frac{2n^2}{3n-1}.$$

    因此
    $$|f(P)-f(Q)|=\left|\frac{n^2}{2n-1}-\frac{2n^2}{3n-1}\right|=\left|n^2\cdot\frac{(3n-1)-(4n-2)}{(2n-1)(3n-1)}\right|=\frac{n^2(n-1)}{|(2n-1)(3n-1)|}.$$

    当 $n\to\infty$ 时，$d(P,Q)\to 0<\delta$，但 $|f(P)-f(Q)|\to+\infty>\varepsilon_0$。故 $f$ 不一致连续。

    Q.E.D.
\end{pf}

\subsection*{Section 9.2}

2. 求下列各函数对于每个自变量的偏微商:

(2) $z = 3^{-\frac{y}{x}}$;
\begin{sol}
    \begin{itemize}
        \item $\dfrac{\partial z}{\partial x}=\ln 3\cdot 3^{-\frac{y}{x}}\cdot\dfrac{y}{x^2}$
        \item $\dfrac{\partial z}{\partial y}=-\ln 3\cdot 3^{-\frac{y}{x}}\cdot\dfrac{1}{x}$
    \end{itemize}
\end{sol}

(4) $z = \ln\left(x + \sqrt{x^2 + y^2}\right)$;
\begin{sol}
    \begin{itemize}
        \item $\dfrac{\partial z}{\partial x}=\dfrac{1+\dfrac{1}{2\sqrt{x^2 + y^2}}\cdot 2x}{x + \sqrt{x^2 + y^2}}=\dfrac{1+\dfrac{x}{\sqrt{x^2 + y^2}}}{x + \sqrt{x^2 + y^2}}=\dfrac{1}{\sqrt{x^2 + y^2}}$
        \item $\dfrac{\partial z}{\partial y}=\dfrac{\dfrac{y}{\sqrt{x^2 + y^2}}}{x + \sqrt{x^2 + y^2}}$
    \end{itemize}
\end{sol}

(6) $u = \mathrm{e}^{x(x^2 + y^2 + z^2)}$;
\begin{sol}
    \begin{itemize}
        \item $\dfrac{\partial u}{\partial x}=\mathrm{e}^{x(x^2 + y^2 + z^2)}(x^2 + y^2 + z^2+2x^2)=\mathrm{e}^{x(x^2 + y^2 + z^2)}(3x^2 + y^2 + z^2)$
        \item $\dfrac{\partial u}{\partial y}=\mathrm{e}^{x(x^2 + y^2 + z^2)}\cdot 2xy$
        \item $\dfrac{\partial u}{\partial z}=\mathrm{e}^{x(x^2 + y^2 + z^2)}\cdot 2xz$
    \end{itemize}
\end{sol}

(8) $u = x\mathrm{e}^{-z} + \ln\left(x + \ln y\right) + z$.
\begin{sol}
    \begin{itemize}
        \item $\dfrac{\partial u}{\partial x}=\mathrm{e}^{-z}+\dfrac{1}{x+\ln y}$
        \item $\dfrac{\partial u}{\partial y}=\dfrac{1}{xy+y\ln y}$
        \item $\dfrac{\partial u}{\partial z}=1-x\mathrm{e}^{-z}$
    \end{itemize}
\end{sol}


3. 设 $f(x,y) = \int_{1}^{x^2 y} \frac{\sin t}{t} \,\mathrm{d}t$, 求 $\dfrac{\partial f}{\partial x}$, $\dfrac{\partial f}{\partial y}$.
\begin{sol}
    \begin{itemize}
        \item $\dfrac{\partial f}{\partial x}=\dfrac{\sin (x^2 y)}{x^2 y}\cdot\dfrac{\partial (x^2y)}{\partial x}=\dfrac{2\sin (x^2 y)}{x}$
        \item $\dfrac{\partial f}{\partial y}=\dfrac{\sin (x^2 y)}{x^2 y}\cdot\dfrac{\partial (x^2y)}{\partial y}=\dfrac{\sin (x^2 y)}{y}$
    \end{itemize}
\end{sol}


4. 设 $f(x,y) =
    \begin{cases}
        y \sin \dfrac{1}{x^2 + y^2}, & x^2 + y^2 \neq 0, \\
        0,                           & x^2 + y^2 = 0.
    \end{cases}$
考察函数 $f(x,y)$ 在原点 $(0,0)$ 的偏导数.
\begin{sol}
    \begin{itemize}
        \item $f_x'(0,0)=\lim\limits_{x\to 0}\dfrac{f(x,0)-f(0,0)}{x}=\lim\limits_{x\to 0}\dfrac{0}{x}=0$
        \item $f_y'(0,0)=\lim\limits_{y\to 0}\dfrac{f(0,y)-f(0,0)}{y}=\lim\limits_{y\to 0}\sin\dfrac{1}{y^2}$ 无限振荡，极限不存在
    \end{itemize}
    
    因此只有 $f_x'(0,0)=0$.
\end{sol}


7. 求曲线
$\begin{cases}
        z = \sqrt{x^2 + y^2 + 1}, \\
        x = 1
    \end{cases}$
上点 $(1,1,\sqrt{3})$ 处的切线分别与 $x$ 轴、$y$ 轴、$z$ 轴正向的夹角.
\begin{sol}
    曲线的参数方程为
    $$\begin{cases}
        x=1 \\
        y=t \\
        z=\sqrt{2+t^2}
    \end{cases},\quad t\in\mathbb{R}.$$
    
    设夹角分别为 $\alpha,\beta,\gamma$.
    
    切线向量为
    $$\boldsymbol{v}=\left.(x',y',z')\right|_{t=1}=\left.\left(0,1,\frac{t}{\sqrt{2+t^2}}\right)\right|_{t=1}=\left(0,1,\frac{1}{\sqrt{3}}\right),$$
    其长度 $|\boldsymbol{v}|=\sqrt{1+\dfrac{1}{3}}=\dfrac{2}{\sqrt{3}}$.
    
    因此
    $$\cos\alpha=0\cdot\frac{\sqrt{3}}{2}=0,$$
    $$\cos\beta=\frac{\sqrt{3}}{2},$$
    $$\cos\gamma=\frac{1}{2}.$$
    
    故 $\alpha=\dfrac{\pi}{2}$, $\beta=\dfrac{\pi}{6}$, $\gamma=\dfrac{\pi}{3}$.
    
    即曲线在 $(1,1,\sqrt{3})$ 处的切线与 $x$ 轴正向的夹角为 $\dfrac{\pi}{2}$，与 $y$ 轴正向的夹角为 $\dfrac{\pi}{6}$，与 $z$ 轴正向的夹角为 $\dfrac{\pi}{3}$.
\end{sol}


8. 证明函数 $u = \dfrac{1}{\sqrt{t}} \mathrm{e}^{-\frac{x^2}{4t}}$ 满足热传导方程 $\dfrac{\partial u}{\partial t} = \dfrac{\partial^2 u}{\partial x^2}$.
\begin{pf}
    首先计算一阶偏导数：
    $$\frac{\partial u}{\partial t}=-\frac{1}{2t^{3/2}}\mathrm{e}^{-\frac{x^2}{4t}}+\frac{1}{\sqrt{t}}\mathrm{e}^{-\frac{x^2}{4t}}\cdot\frac{x^2}{4t^2}=\frac{x^2-2t}{4t^{5/2}}\mathrm{e}^{-\frac{x^2}{4t}},$$
    $$\frac{\partial u}{\partial x}=\frac{1}{\sqrt{t}} \mathrm{e}^{-\frac{x^2}{4t}}\cdot\left(-\frac{x}{2t}\right).$$
    
    再计算二阶偏导数：
    $$\frac{\partial^2 u}{\partial x^2}=\frac{1}{\sqrt{t}} \mathrm{e}^{-\frac{x^2}{4t}}\cdot\left(-\frac{x^2}{4t^2}\right)-\frac{1}{2t^{3/2}} \mathrm{e}^{-\frac{x^2}{4t}}=\frac{x^2-2t}{4t^{5/2}}\mathrm{e}^{-\frac{x^2}{4t}}.$$
    
    因此 $\dfrac{\partial u}{\partial t} = \dfrac{\partial^2 u}{\partial x^2}$.
    
    Q.E.D.
\end{pf}

\end{document}